% ===== PRÉAMBULE CANONIQUE — à coller VERBATIM en tête de CHAQUE .tex =====
\documentclass[11pt,a4paper]{article}
\usepackage[utf8]{inputenc}\usepackage[T1]{fontenc}\usepackage{lmodern}
\usepackage[french]{babel}
\usepackage{amsmath,amssymb,mathtools}\usepackage{icomma}
\usepackage{siunitx}\sisetup{locale=FR}  % \qty{}{}, \si{}, \ang{}
\usepackage{xcolor}\usepackage{colortbl}
\usepackage{tikz}\usepackage{pgfplots}\pgfplotsset{compat=1.18}
\usepackage[siunitx,europeanresistors]{circuitikz} % circuits (élec.)
\usepackage[most]{tcolorbox}\usepackage{enumitem}\usepackage{array,booktabs}
\usepackage[a4paper,margin=2.2cm]{geometry}
\usetikzlibrary{arrows.meta,calc,angles,quotes,positioning,patterns,%
  backgrounds,shapes.geometric,decorations.pathmorphing,decorations.markings,intersections}
\definecolor{primary}{RGB}{222,49,99}\definecolor{primarydark}{RGB}{150,22,55}
\definecolor{defcol}{RGB}{0,114,178}\definecolor{methcol}{RGB}{52,92,148}
\definecolor{excol}{RGB}{0,135,90}\definecolor{attncol}{RGB}{200,40,55}
\tcbset{boxbase/.style={enhanced,breakable,boxrule=0.9pt,arc=2.5pt,
  left=3mm,right=3mm,top=2.3mm,bottom=2.3mm,fonttitle=\bfseries,coltitle=white}}
% --- boîtes TUTORIEL (noms EXACTS, ne pas renommer) ---
\newtcolorbox{cle}[1][]{boxbase,colback=defcol!6,colframe=defcol,
  title={\textsc{À retenir}\ifx\relax#1\relax\relax\else\textmd{ : \itshape #1}\fi}}
\newtcolorbox{methode}[1][]{boxbase,colback=methcol!7,colframe=methcol,sharp corners,
  title={\textsc{Méthode}\ifx\relax#1\relax\relax\else\textmd{ --- \itshape #1}\fi}}
\newtcolorbox{exemplebox}[1][]{boxbase,colback=excol!5,colframe=excol,
  title={\textsc{Exemple}\ifx\relax#1\relax\relax\else\textmd{ : \itshape #1}\fi}}
\newtcolorbox{piege}[1][]{boxbase,colback=attncol!6,colframe=attncol,
  title={\textsc{Piège}\ifx\relax#1\relax\relax\else\textmd{ : \itshape #1}\fi}}
\newtcolorbox{remarque}[1][]{boxbase,colback=black!4,colframe=black!45,
  title={\textsc{Remarque}\ifx\relax#1\relax\relax\else\textmd{ : \itshape #1}\fi}}
% --- boîtes EXERCICE / CORRIGÉ (noms EXACTS, auto-comptées) ---
\newtcolorbox[auto counter]{exo}[1][]{boxbase,colback=primary!4,colframe=primary,
  title={\textsc{Exercice}~\thetcbcounter\ifx\relax#1\relax\relax\else\textmd{ --- \itshape #1}\fi}}
\newtcolorbox[auto counter]{corrige}[1][]{boxbase,colback=excol!4,colframe=excol,
  title={\textsc{Corrigé}~\thetcbcounter\ifx\relax#1\relax\relax\else\textmd{ --- \itshape #1}\fi}}
\newcommand{\vv}[1]{\vec{#1}}\newcommand{\ds}{\displaystyle}
\newcommand{\R}{\mathbb{R}}\newcommand{\N}{\mathbb{N}}\newcommand{\Z}{\mathbb{Z}}
% (un \usepackage{fancyhdr}+en-tête est possible ; il est ignoré côté web)
% =========================================================================
\begin{document}
% THEME_TITLE: La quantité de mouvement
\sisetup{per-mode=symbol}

À côté de l'énergie, une autre grandeur mesure « l'élan » d'un corps en mouvement : la
\textbf{quantité de mouvement}. Contrairement à l'énergie (scalaire), c'est une grandeur
\emph{vectorielle}.

\begin{cle}[quantité de mouvement]
\[
\boxed{\;\vv p=m\,\vv v\;}\qquad(\text{en }\si{\kilogram\metre\per\second})
\]
$m$ en \si{\kilogram}, $\vv v$ en \si{\metre\per\second}. Le vecteur $\vv p$ a la \emph{même direction}
et le \emph{même sens} que la vitesse $\vv v$. En norme : $p=m\,v$. \textbf{Attention} : $v$ doit être
en \si{\metre\per\second} (diviser les \si{\kilo\metre\per\hour} par $3,6$).
\end{cle}

\begin{center}
\begin{tikzpicture}[>=Stealth]
  % voiture
  \draw[fill=defcol!15,draw=black!70,rounded corners=2pt] (0,0) rectangle (1.2,0.5);
  \node at (0.6,0.25){\scriptsize voiture};
  \draw[->,line width=1.2pt,defcol] (1.3,0.25)--(3.3,0.25) node[midway,above]{\scriptsize $\vv p_1=m_1\vv v$};
  % camion
  \draw[fill=primary!15,draw=black!70,rounded corners=2pt] (0,-1.4) rectangle (1.7,-0.8);
  \node at (0.85,-1.1){\scriptsize camion};
  \draw[->,line width=1.2pt,primary] (1.8,-1.1)--(5.8,-1.1) node[midway,above]{\scriptsize $\vv p_2=m_2\vv v$};
  \node[black!60,align=left] at (7.7,-0.4){\scriptsize À la \emph{même} vitesse $\vv v$,\\ \scriptsize le camion ($m_2=2m_1$) a\\ \scriptsize une quantité de mouvement\\ \scriptsize \emph{deux fois} plus grande.};
\end{tikzpicture}
\end{center}

\begin{methode}[calculer et comparer des quantités de mouvement]
\begin{enumerate}[leftmargin=1.6em,topsep=1pt,itemsep=1pt]
  \item Convertir la vitesse en \si{\metre\per\second} si elle est donnée en \si{\kilo\metre\per\hour}.
  \item Calculer $p=m\,v$ pour chaque mobile.
  \item Pour comparer, faire le \emph{rapport} des quantités de mouvement (à même vitesse,
        $\dfrac{p_2}{p_1}=\dfrac{m_2}{m_1}$).
\end{enumerate}
\end{methode}

\begin{cle}[quantité de mouvement et « difficulté à arrêter »]
Plus la quantité de mouvement d'un corps est grande, plus il est \emph{difficile à arrêter} : il faut
une force plus grande, ou une distance (ou durée) de freinage plus longue. Un corps \emph{massif et
rapide} a une grande quantité de mouvement.
\end{cle}

\begin{remarque}[ne pas confondre $\vv p$ et $E_c$]
Les deux augmentent avec la masse et la vitesse, mais différemment :
\[
p=m\,v\quad(\text{proportionnel à }v,\ \text{vectoriel}),\qquad
E_c=\tfrac12 m v^2\quad(\text{en }v^2,\ \text{scalaire}).
\]
Doubler la vitesse \emph{double} $p$ mais \emph{quadruple} $E_c$. On peut relier les deux :
$E_c=\dfrac{p^2}{2m}$.
\end{remarque}

\begin{piege}[erreurs fréquentes]
\begin{itemize}[leftmargin=1.5em,topsep=1pt,itemsep=1pt]
  \item Donner $p$ en \si{\joule} : l'unité est le \si{\kilogram\metre\per\second}, pas le joule.
  \item Oublier que $\vv p$ est un \textbf{vecteur} (direction et sens de $\vv v$).
  \item Utiliser une vitesse en \si{\kilo\metre\per\hour} sans la convertir.
\end{itemize}
\end{piege}
\end{document}
